% Exemplo de uso do preâmbulo e comandos personalizados

\documentclass[a4paper,landscape]{article}
\usepackage{latexconfig/preambulo_formularios}

\begin{document}

\begin{multicols}{3}

\formulabox[blue]{F = m a}{F = Força (N) \\ m = Massa (kg) \\ a = Aceleração (m/s$^2$)}

\formulabox[red]{E = m c$^2$}{E = Energia (J) \\ m = Massa (kg) \\ c = Velocidade da luz ($3{,}0 \times 10^8$ m/s)}

\formulabox{P = F / A}{P = Pressão (Pa) \\ F = Força (N) \\ A = Área (m$^2$)}

\formulabox[teal]{Q = m c \Delta T}{Q = Calor (J) \\ m = Massa (kg) \\ c = Calor específico (J/kg$\cdot$K) \\ $\Delta T$ = Variação de temperatura (K)}

\formulabox[orange]{F = k \frac{q_1 q_2}{r^2}}{F = Força elétrica (N) \\ k = 9,0 $\times$ 10$^9$ N m$^2$/C$^2$ \\ $q_1$, $q_2$ = Cargas (C) \\ r = Distância (m)}

\textbf{Constantes úteis:} \\
\constbox[blue]{\pi}{3,1415}
\constbox[red]{e}{2,718}
\constbox{g}{9,8 m/s$^2$}
\constbox[teal]{c}{3,0 $\times$ 10$^8$ m/s}
\constbox[orange]{\epsilon_0}{8,85 $\times$ 10$^{-12}$ C$^2$/N m$^2$}
\constbox[khaki]{k_B}{1,38 $\times$ 10$^{-23}$ J/K}

\end{multicols}

\end{document}
